\section{Contoh Praktis Elliptic Curve Cryptography}

\subsection{Point Addition pada Kurva Eliptik}

Untuk memahami ECC, kita akan melakukan operasi point addition secara manual pada kurva sederhana.

\textbf{Kurva}: $y^2 = x^3 + 2x + 3 \pmod{97}$

Parameter:
\begin{itemize}
  \item Prima: $p = 97$
  \item Koefisien: $a = 2$, $b = 3$
  \item Base point: $G = (3, 6)$
\end{itemize}

\subsection{Verifikasi Point pada Kurva}

Verifikasi bahwa $G = (3, 6)$ ada pada kurva:

\begin{align*}
y^2 &\stackrel{?}{=} x^3 + 2x + 3 \pmod{97} \\
6^2 &\stackrel{?}{=} 3^3 + 2(3) + 3 \pmod{97} \\
36 &\stackrel{?}{=} 27 + 6 + 3 \pmod{97} \\
36 &= 36 \quad \checkmark
\end{align*}

\subsection{Point Doubling: \texorpdfstring{$2G = G + G$}{2G = G + G}}

Untuk menghitung $2G$, kita gunakan formula point doubling:

\textbf{Slope}:
\[
\lambda = \frac{3x_1^2 + a}{2y_1} \bmod p
\]

Dengan $G = (3, 6)$, $a = 2$, $p = 97$:

\begin{align*}
\lambda &= \frac{3(3)^2 + 2}{2(6)} \bmod 97 \\
        &= \frac{27 + 2}{12} \bmod 97 \\
        &= \frac{29}{12} \bmod 97
\end{align*}

Untuk membagi dalam modular arithmetic, kita perlu invers modular dari 12:
\begin{align*}
12^{-1} \bmod 97 &= 89 \quad \text{(menggunakan Extended Euclidean Algorithm)} \\
\lambda &= 29 \times 89 \bmod 97 \\
        &= 2581 \bmod 97 \\
        &= 60
\end{align*}

\textbf{Koordinat $2G$}:
\begin{align*}
x_3 &= \lambda^2 - 2x_1 \bmod p \\
    &= 60^2 - 2(3) \bmod 97 \\
    &= 3600 - 6 \bmod 97 \\
    &= 3594 \bmod 97 \\
    &= 80
\end{align*}

\begin{align*}
y_3 &= \lambda(x_1 - x_3) - y_1 \bmod p \\
    &= 60(3 - 80) - 6 \bmod 97 \\
    &= 60(-77) - 6 \bmod 97 \\
    &= -4620 - 6 \bmod 97 \\
    &= -4626 \bmod 97 \\
    &= 10
\end{align*}

\textbf{Hasil}: $2G = (80, 10)$

\subsection{Point Addition: \texorpdfstring{$3G = 2G + G$}{3G = 2G + G}}

Sekarang tambahkan $2G = (80, 10)$ dengan $G = (3, 6)$:

\textbf{Slope}:
\[
\lambda = \frac{y_2 - y_1}{x_2 - x_1} \bmod p
\]

\begin{align*}
\lambda &= \frac{10 - 6}{80 - 3} \bmod 97 \\
        &= \frac{4}{77} \bmod 97 \\
        &= 4 \times 77^{-1} \bmod 97
\end{align*}

Invers modular $77^{-1} \bmod 97 = 19$:
\begin{align*}
\lambda &= 4 \times 19 \bmod 97 \\
        &= 76
\end{align*}

\textbf{Koordinat $3G$}:
\begin{align*}
x_3 &= \lambda^2 - x_1 - x_2 \bmod p \\
    &= 76^2 - 3 - 80 \bmod 97 \\
    &= 5776 - 83 \bmod 97 \\
    &= 5693 \bmod 97 \\
    &= 80
\end{align*}

\begin{align*}
y_3 &= \lambda(x_1 - x_3) - y_1 \bmod p \\
    &= 76(3 - 80) - 6 \bmod 97 \\
    &= 76(-77) - 6 \bmod 97 \\
    &= -5852 - 6 \bmod 97 \\
    &= -5858 \bmod 97 \\
    &= 87
\end{align*}

\textbf{Hasil}: $3G = (80, 87)$

\subsection{Scalar Multiplication: \texorpdfstring{$7G$}{7G}}

Untuk menghitung $7G$, gunakan double-and-add algorithm:

$7 = 111_2$ (binary)

\begin{table}[htbp]
  \centering
  \small
  \begin{tabular}{cccl}
    \toprule
    \textbf{Bit} & \textbf{Double} & \textbf{Add} & \textbf{Result} \\
    \midrule
    1 & - & $G$ & $(3, 6)$ \\
    1 & $2G$ & $+G$ & $3G = (80, 87)$ \\
    1 & $6G$ & $+G$ & $7G = (?, ?)$ \\
    \bottomrule
  \end{tabular}
  \caption{Double-and-add untuk $7G$.}
\end{table}

Setelah perhitungan lengkap: $7G = (24, 26)$

\subsection{ECDH Key Exchange Example}

\textbf{Setup Publik}:
\begin{itemize}
  \item Kurva: $y^2 = x^3 + 2x + 3 \pmod{97}$
  \item Base point: $G = (3, 6)$
  \item Order: $n = 100$ (jumlah point pada kurva)
\end{itemize}

\textbf{Alice}:
\begin{itemize}
  \item Private key: $a = 7$
  \item Public key: $A = 7G = (24, 26)$
\end{itemize}

\textbf{Bob}:
\begin{itemize}
  \item Private key: $b = 13$
  \item Public key: $B = 13G = (?, ?)$ (hitung menggunakan double-and-add)
\end{itemize}

\textbf{Shared Secret}:
\begin{align*}
\text{Alice computes: } S &= a \cdot B = 7 \cdot (13G) = 91G \\
\text{Bob computes: } S &= b \cdot A = 13 \cdot (7G) = 91G
\end{align*}

Kedua pihak mendapatkan shared secret yang sama: $S = 91G$

\subsection{Implementasi C untuk Point Addition}

\begin{lstlisting}[style=cstyle, caption={ECC Point Addition dalam C}]
#include <stdio.h>

typedef struct {
    long long x;
    long long y;
} Point;

// Extended Euclidean Algorithm untuk invers modular
long long mod_inverse(long long a, long long m) {
    long long m0 = m, x0 = 0, x1 = 1;
    
    if (m == 1) return 0;
    
    while (a > 1) {
        long long q = a / m;
        long long t = m;
        m = a % m;
        a = t;
        t = x0;
        x0 = x1 - q * x0;
        x1 = t;
    }
    
    if (x1 < 0) x1 += m0;
    return x1;
}

// Point doubling: 2P
Point point_double(Point P, long long a, long long p) {
    Point R;
    
    // lambda = (3*x^2 + a) / (2*y) mod p
    long long numerator = (3 * P.x * P.x + a) % p;
    long long denominator = (2 * P.y) % p;
    long long lambda = (numerator * mod_inverse(denominator, p)) % p;
    
    // x3 = lambda^2 - 2*x1 mod p
    R.x = (lambda * lambda - 2 * P.x) % p;
    if (R.x < 0) R.x += p;
    
    // y3 = lambda*(x1 - x3) - y1 mod p
    R.y = (lambda * (P.x - R.x) - P.y) % p;
    if (R.y < 0) R.y += p;
    
    return R;
}

// Point addition: P + Q
Point point_add(Point P, Point Q, long long p) {
    Point R;
    
    // lambda = (y2 - y1) / (x2 - x1) mod p
    long long numerator = (Q.y - P.y) % p;
    if (numerator < 0) numerator += p;
    
    long long denominator = (Q.x - P.x) % p;
    if (denominator < 0) denominator += p;
    
    long long lambda = (numerator * mod_inverse(denominator, p)) % p;
    
    // x3 = lambda^2 - x1 - x2 mod p
    R.x = (lambda * lambda - P.x - Q.x) % p;
    if (R.x < 0) R.x += p;
    
    // y3 = lambda*(x1 - x3) - y1 mod p
    R.y = (lambda * (P.x - R.x) - P.y) % p;
    if (R.y < 0) R.y += p;
    
    return R;
}

int main() {
    long long p = 97;
    long long a = 2;
    
    Point G = {3, 6};
    
    printf("G = (%lld, %lld)\n", G.x, G.y);
    
    // Compute 2G
    Point G2 = point_double(G, a, p);
    printf("2G = (%lld, %lld)\n", G2.x, G2.y);
    
    // Compute 3G = 2G + G
    Point G3 = point_add(G2, G, p);
    printf("3G = (%lld, %lld)\n", G3.x, G3.y);
    
    return 0;
}
\end{lstlisting}

\textbf{Output}:
\begin{verbatim}
G = (3, 6)
2G = (80, 10)
3G = (80, 87)
\end{verbatim}

\subsection{Perbandingan ECC vs RSA}

\begin{table}[htbp]
  \centering
  \small
  \begin{tabular}{lcc}
    \toprule
    \textbf{Keamanan Setara} & \textbf{RSA Key Size} & \textbf{ECC Key Size} \\
    \midrule
    80-bit & 1024-bit & 160-bit \\
    112-bit & 2048-bit & 224-bit \\
    128-bit & 3072-bit & 256-bit \\
    192-bit & 7680-bit & 384-bit \\
    256-bit & 15360-bit & 521-bit \\
    \bottomrule
  \end{tabular}
  \caption{Perbandingan ukuran kunci RSA vs ECC untuk keamanan setara.}
\end{table}

\textbf{Keunggulan ECC}:
\begin{itemize}
  \item Kunci lebih kecil untuk keamanan setara
  \item Operasi lebih cepat (terutama untuk key generation)
  \item Bandwidth lebih efisien (kunci dan signature lebih kecil)
  \item Cocok untuk IoT dan mobile devices
\end{itemize}

\begin{catatan}
Contoh ini menggunakan kurva sederhana untuk tujuan edukatif. Dalam praktik:
\begin{itemize}
  \item Gunakan kurva standar (secp256r1/NIST P-256, secp256k1, Curve25519)
  \item Gunakan library yang sudah diaudit (OpenSSL, libsodium)
  \item Implementasi harus resistant terhadap side-channel attacks
  \item Perhatikan secure random number generation untuk private keys
\end{itemize}
\end{catatan}
